\documentclass[12pt,a4paper,oneside]{report}
\usepackage[utf8]{inputenc}
\usepackage[T1]{fontenc}
\usepackage[brazil]{babel}
\usepackage{lmodern,microtype}
\usepackage[a4paper,top=2.7cm,bottom=2.5cm,left=3cm,right=2.2cm]{geometry}
\usepackage{setspace,graphicx,booktabs,longtable,array,tabularx,enumitem,xcolor,listings,titlesec,hyperref,fancyhdr,float,amsmath,caption}
\definecolor{melodixblue}{HTML}{244A73}
\definecolor{codebg}{HTML}{F4F6F8}
\definecolor{codeframe}{HTML}{D4DCE4}
\definecolor{muted}{HTML}{536273}
\hypersetup{colorlinks=true,linkcolor=melodixblue,urlcolor=melodixblue,citecolor=melodixblue,pdftitle={Documentação Técnica do Melodix},pdfsubject={Arquitetura, APIs, segurança, persistência e áudio}}
\setstretch{1.18}
\setlength{\parindent}{1.25cm}
\setlength{\parskip}{0.25em}
\setlist[itemize]{leftmargin=1.2cm,itemsep=0.2em}
\setlist[enumerate]{leftmargin=1.2cm,itemsep=0.2em}
\renewcommand{\arraystretch}{1.2}
\titleformat{\chapter}[display]{\normalfont\bfseries\color{melodixblue}}{\large\chaptername\ \thechapter}{0.5em}{\Huge}
\titleformat{\section}{\normalfont\Large\bfseries\color{melodixblue}}{\thesection}{0.7em}{}
\titleformat{\subsection}{\normalfont\large\bfseries}{\thesubsection}{0.7em}{}
\titleformat{\subsubsection}{\normalfont\normalsize\bfseries}{\thesubsubsection}{0.7em}{}
\pagestyle{fancy}
\fancyhf{}
\fancyhead[L]{\small Documentação Técnica do Melodix}
\fancyhead[R]{\small Projeto de streaming musical}
\fancyfoot[C]{\thepage}
\renewcommand{\headrulewidth}{0.4pt}
\lstdefinestyle{codigo}{basicstyle=\ttfamily\small,backgroundcolor=\color{codebg},frame=single,rulecolor=\color{codeframe},breaklines=true,columns=fullflexible,keepspaces=true,showstringspaces=false,tabsize=2,numbers=left,numberstyle=\tiny\color{muted},xleftmargin=1.5em,framexleftmargin=1em}
\lstset{style=codigo}
\newcommand{\arquivo}[1]{\texttt{\detokenize{#1}}}
\begin{document}
\begin{titlepage}
\centering
\vspace*{2.2cm}
{\Large DOCUMENTAÇÃO TÉCNICA\par}
\vspace{1.1cm}
{\Huge\bfseries\color{melodixblue} Melodix\par}
\vspace{0.6cm}
{\LARGE Aplicação de streaming musical\par}
\vspace{1.8cm}
\rule{0.78\textwidth}{0.6pt}\par
\vspace{0.7cm}
{\large Arquitetura, tecnologias, serviços, APIs, segurança, persistência e processamento de áudio\par}
\vspace{0.7cm}
\rule{0.78\textwidth}{0.6pt}\par
\vfill
{\large Documento elaborado a partir da inspeção do código-fonte, arquivos de configuração, rotas e migrações SQL.\par}
\vspace{1.4cm}
{\large Outubro de 2026\par}
\end{titlepage}
\pagenumbering{roman}
\tableofcontents
\listoftables
\clearpage
\pagenumbering{arabic}

\chapter*{Apresentação}
\addcontentsline{toc}{chapter}{Apresentação}
Este documento descreve a implementação técnica do Melodix, aplicação de reprodução e organização de músicas com interface web e possibilidade de empacotamento para Android. A análise considera o README, arquivos \arquivo{package.json}, serviços Angular, módulos Express, configurações do Supabase e migrações SQL.

O Melodix utiliza o catálogo público do Audius para localizar faixas e obter endereços de reprodução. A autenticação e a biblioteca pessoal são apoiadas pelo Supabase. O navegador reproduz o áudio por meio de um elemento HTML de áudio e, quando as condições de origem permitem, disponibiliza amostras espectrais para o visualizador construído com a Web Audio API.

As versões indicadas são as declarações de dependência encontradas nos manifestos do projeto. Faixas com \texttt{\string^} ou \texttt{\string~} representam intervalos semver, não uma confirmação da versão efetivamente instalada. O README informa Node.js 18 ou superior.

\chapter{Visão geral do sistema}
\section{Objetivo e funcionalidades}
O Melodix reúne pesquisa e reprodução musical com recursos de biblioteca individual. O README descreve:
\begin{itemize}
\item cadastro e autenticação por e-mail e senha via Supabase Auth;
\item pesquisa de faixas no catálogo Audius;
\item paginação dos resultados e carregamento adicional durante a rolagem;
\item reprodução, pausa, faixa anterior e próxima, progresso e volume;
\item visualizador de áudio com Web Audio API e alternativa animada em CSS;
\item criação e manutenção de playlists e músicas favoritas;
\item execução em navegador e empacotamento Android por Capacitor.
\end{itemize}

\section{Limites funcionais observados}
O catálogo e a origem do áudio são providos pelo Audius. O projeto não mantém um servidor próprio para armazenar ou distribuir os arquivos musicais. O backend Melodix intermedeia consultas de catálogo e mantém os dados da biblioteca, enquanto o cliente de áudio acessa diretamente o endereço de streaming.

Existe uma migração SQL denominada \arquivo{20260518002513\_spotify\_connections.sql}, que cria uma tabela para conexões Spotify. A existência dessa tabela, isoladamente, não comprova um fluxo OAuth funcional. Na inspeção realizada, não foi identificada implementação correspondente de autorização, callback e renovação de tokens Spotify. A integração de catálogo e reprodução efetivamente identificada está baseada no Audius.

\section{Organização em camadas}
A solução divide-se em três áreas:
\begin{enumerate}
\item \textbf{Frontend}: Angular e Ionic, responsáveis por telas, interação, estado de reprodução e chamadas HTTP.
\item \textbf{Backend}: API HTTP Express e TypeScript, responsável por validar requisições, autenticar operações da biblioteca, consultar Audius e acessar Supabase.
\item \textbf{Serviços externos e persistência}: Audius para catálogo e stream, Supabase Auth para identidade e PostgreSQL para dados da biblioteca.
\end{enumerate}

\chapter{Tecnologias, versões e ferramentas}
\section{Frontend}
A aplicação cliente utiliza Angular 20, Ionic 8 e TypeScript. Angular organiza componentes, serviços, rotas e interceptors; Ionic fornece componentes visuais adaptados a aplicações web e móveis.

\begin{longtable}{p{0.30\textwidth}p{0.22\textwidth}p{0.40\textwidth}}
\toprule
\textbf{Dependência} & \textbf{Versão declarada} & \textbf{Responsabilidade}\\
\midrule
Angular & \texttt{\string^20.0.0} & Framework da aplicação cliente.\\
Angular CLI & \texttt{20.3.26} & Compilação e ferramentas de desenvolvimento.\\
Ionic Angular & \texttt{\string^8.0.0} & Componentes de interface.\\
RxJS & \texttt{\string~7.8.0} & Fluxos assíncronos e estado observável.\\
TypeScript & \texttt{\string~5.9.0} & Tipagem e compilação do frontend.\\
Capacitor Core & \texttt{8.3.4} & Ponte entre aplicação web e recursos nativos.\\
Capacitor CLI & \texttt{8.3.4} & Sincronização do projeto nativo.\\
Supabase JS & \texttt{\string^2.105.4} & Cliente JavaScript do Supabase.\\
Ionicons & \texttt{\string^7.0.0} & Ícones da interface.\\
\bottomrule
\end{longtable}
O manifesto também declara Jasmine, Karma, ESLint e ferramentas Angular ESLint para testes e qualidade.

\section{Backend}
O backend é executado em Node.js, utiliza Express 4 e é escrito em TypeScript, compilado para JavaScript CommonJS.

\begin{longtable}{p{0.30\textwidth}p{0.22\textwidth}p{0.40\textwidth}}
\toprule
\textbf{Dependência} & \textbf{Versão declarada} & \textbf{Responsabilidade}\\
\midrule
Express & \texttt{\string^4.21.2} & Servidor HTTP e roteamento.\\
TypeScript & \texttt{\string^5.7.2} & Tipagem e compilação.\\
tsx & \texttt{\string^4.19.2} & Execução de TypeScript em desenvolvimento.\\
Axios & \texttt{\string^1.7.9} & Requisições HTTP ao Audius.\\
Supabase JS & \texttt{\string^2.105.4} & Integração com Auth e PostgreSQL.\\
Zod & \texttt{\string^3.24.1} & Validação de variáveis de ambiente.\\
cors & \texttt{\string^2.8.5} & Política de compartilhamento entre origens.\\
dotenv & \texttt{\string^16.4.7} & Carregamento de variáveis de ambiente.\\
\bottomrule
\end{longtable}
O \arquivo{backend/tsconfig.json} define ES2022, CommonJS, resolução Node, modo estrito e saída em \arquivo{dist}. O README informa Node.js 18 ou superior.

\section{Ambiente e publicação}
Para desenvolvimento, o frontend é iniciado com \texttt{ionic serve}; o backend pode ser iniciado por \texttt{npm run backend:dev}. A compilação do backend utiliza \texttt{npm run build} e a execução de produção utiliza \texttt{node dist/server.js}.

O README registra hospedagem do backend no Render, com diretório raiz \arquivo{backend}, comando de build \texttt{npm install \&\& npm run build} e comando de inicialização \texttt{npm start}. Esses dados representam a configuração documentada no repositório, não uma verificação independente do estado atual do serviço remoto.

\section{Aplicativo Android}
O README descreve a geração do aplicativo Android pelo Capacitor. O processo consiste em compilar a aplicação Ionic em produção, sincronizar os arquivos nativos com \texttt{npx cap sync android} e gerar APK de depuração pelo Gradle. A configuração efetiva do SDK depende do ambiente Android Studio instalado.

\chapter{Arquitetura do frontend}
\section{Estrutura de diretórios}
A organização principal identificada é:
\begin{lstlisting}
src/
  app/
    pages/
      home/ search/ library/ playlist/ player/ login/
    components/
      player/ audio-visualizer/ playlists/
    core/
      services/ interceptors/ guards/
    shared/
      components/ models/ directives/ pipes/
  environments/
\end{lstlisting}
As páginas representam contextos de navegação. Os componentes encapsulam partes reutilizáveis da interface. A pasta \arquivo{core} concentra serviços e mecanismos transversais. A pasta \arquivo{shared} reúne modelos, diretivas, pipes e componentes compartilhados.

\section{CoreModule e injeção de dependências}
O arquivo \arquivo{src/app/core/core.module.ts} registra \texttt{SupabaseService}, \texttt{AuthService}, \texttt{AudiusService}, \texttt{PlayerService} e \texttt{PlaylistService}. Registra também \texttt{ApiInterceptor} no token \texttt{HTTP\_INTERCEPTORS}, com \texttt{multi: true}.

O construtor utiliza \texttt{@Optional()} e \texttt{@SkipSelf()} para detectar uma importação anterior do módulo. Se já houver instância-pai, lança exceção. A verificação evita registros duplicados dos serviços centrais.

\section{Serviços principais}
\subsection{AudiusService}
O serviço Angular \arquivo{src/app/core/services/audius.service.ts} utiliza \texttt{HttpClient}. O método \texttt{searchTracks(query, limit, offset)} monta parâmetros HTTP e consulta \texttt{/audius/search}. A resposta é transformada em \texttt{tracks} e \texttt{hasMore}. O método \texttt{getTrack(id)} consulta \texttt{/audius/track/:id}. Ambos retornam \texttt{Observable}, permitindo composição com operadores RxJS.

\subsection{PlayerService}
O serviço \arquivo{src/app/core/services/player.service.ts} centraliza a reprodução. Mantém uma instância de \texttt{Audio} e publica estados por \texttt{BehaviorSubject} e \texttt{Subject}. Entre os estados observáveis estão faixa atual, reprodução, progresso, carregamento, volume, mensagens e nó analisador.

\texttt{BehaviorSubject} entrega o último valor conhecido a novos assinantes. \texttt{Subject} é usado para eventos que não precisam reter estado, como solicitações de faixa anterior ou próxima. O serviço escuta \texttt{timeupdate}, \texttt{ended}, \texttt{waiting} e \texttt{playing} para atualizar os estados da interface.

\subsection{PlaylistService e AuthService}
O \texttt{PlaylistService} concentra operações da biblioteca e favoritos. O \texttt{AuthService} concentra o fluxo de autenticação do cliente com Supabase. O interceptor HTTP aplica comportamento transversal às chamadas. A implementação exata de cada cabeçalho deve ser conferida diretamente no interceptor em futuras alterações.

\section{Componentes de interface}
O \texttt{PlayerComponent} oferece play/pause, navegação, volume, progresso, favoritos e inclusão em playlists. A propriedade \texttt{embedded} permite apresentação compacta na tela de busca.

O componente raiz \texttt{AppComponent} combina a faixa atual com eventos de navegação. O mini-player aparece quando existe faixa selecionada e a rota não começa com \texttt{/search}. Essa regra utiliza \texttt{combineLatest}, \texttt{map}, \texttt{filter} e \texttt{startWith}.

A tela de busca apresenta a roda de seleção, lista de resultados e carregamento de páginas adicionais. O template usa \texttt{trackBy} para identificar itens pela identidade da faixa e reduzir atualizações desnecessárias.

\section{Ciclo de vida e inscrições}
No componente raiz, \texttt{destroy\$} é utilizado com \texttt{takeUntil} para encerrar inscrições quando o componente é destruído. O método \texttt{ngOnDestroy} emite e conclui o Subject.

O visualizador encerra sua inscrição RxJS, desconecta o \texttt{ResizeObserver} e cancela o ciclo de animação com \texttt{cancelAnimationFrame}. Isso evita observadores ativos após a remoção do componente.

\chapter{Arquitetura do backend}
\section{Estrutura de diretórios}
\begin{lstlisting}
backend/
  src/
    server.ts
    config/env.ts
    lib/supabase.ts
    middlewares/error.middleware.ts
    modules/
      auth/auth.middleware.ts
      audius/
        audius-http.ts
        audius.controller.ts
        audius.routes.ts
        audius.service.ts
      library/
        library.controller.ts
        library.routes.ts
        library.schemas.ts
        library.service.ts
    types/track.ts
    utils/async-handler.ts
    utils/http-error.ts
\end{lstlisting}
A divisão modular separa transporte HTTP, regras de negócio, integração externa, autenticação, configuração e utilitários.

\section{Inicialização do servidor}
O arquivo \arquivo{backend/src/server.ts} cria a aplicação Express e lê porta e URL do frontend a partir da configuração validada. Configura CORS, interpreta JSON com limite de 1 MB, registra \texttt{/health}, monta os roteadores \texttt{/library} e \texttt{/audius} e instala o middleware central de erro.

A ordem é relevante: CORS e parser JSON são instalados antes dos roteadores; o middleware de erro é registrado depois deles para receber falhas encaminhadas pelas rotas e controladores.

\section{Controladores, serviços e utilitários}
As rotas associam métodos HTTP a controladores. Os controladores interpretam parâmetros, corpo e identidade do usuário, delegando operações aos serviços. Os serviços concentram chamadas ao Audius ou Supabase.

O utilitário \texttt{asyncHandler} encaminha rejeições de funções assíncronas ao fluxo de erro do Express. \texttt{HttpError} representa falhas HTTP com mensagem e status. O middleware centraliza as respostas de erro.

\section{Configuração por ambiente}
O arquivo \arquivo{backend/src/config/env.ts} executa \texttt{dotenv.config()} e valida \texttt{process.env} com Zod. São definidos:
\begin{itemize}
\item \texttt{PORT}: inteiro positivo, padrão 3000;
\item \texttt{FRONTEND\_URL}: URL válida, padrão \texttt{http://localhost:8100};
\item \texttt{SUPABASE\_URL}: URL obrigatória;
\item \texttt{SUPABASE\_PUBLISHABLE\_KEY}: chave pública obrigatória;
\item \texttt{SUPABASE\_SERVICE\_ROLE\_KEY}: chave de serviço obrigatória;
\item \texttt{AUDIUS\_APP\_NAME}: nome do app, padrão Melodix.
\end{itemize}
A chamada \texttt{envSchema.parse(process.env)} interrompe a inicialização quando variáveis obrigatórias não atendem ao esquema, antecipando falhas de configuração.

\chapter{APIs e integrações externas}
\section{API interna}
As rotas são agrupadas por domínio. O roteador Audius expõe consulta pública do catálogo. O roteador Library aplica autenticação antes de registrar operações.

\begin{longtable}{p{0.13\textwidth}p{0.48\textwidth}p{0.33\textwidth}}
\toprule
\textbf{Método} & \textbf{Rota} & \textbf{Descrição}\\
\midrule
GET & \texttt{/health} & Estado básico do backend.\\
GET & \texttt{/audius/search} & Pesquisa de faixas.\\
GET & \texttt{/audius/track/:id} & Consulta de faixa por identificador.\\
GET & \texttt{/library/playlists} & Lista playlists do usuário.\\
POST & \texttt{/library/playlists} & Cria playlist.\\
GET & \texttt{/library/playlists/:id} & Consulta playlist e faixas.\\
PUT & \texttt{/library/playlists/:id} & Atualiza playlist.\\
DELETE & \texttt{/library/playlists/:id} & Exclui playlist.\\
POST & \texttt{/library/playlists/:id/tracks} & Adiciona ou atualiza faixa.\\
DELETE & \texttt{/library/playlists/:id/tracks/:trackId} & Remove faixa.\\
GET & \texttt{/library/favorites} & Lista favoritos.\\
POST & \texttt{/library/favorites} & Registra favorito.\\
DELETE & \texttt{/library/favorites/:trackId} & Remove favorito.\\
\bottomrule
\end{longtable}
As rotas \texttt{/library} são protegidas por \texttt{authMiddleware}. As rotas \texttt{/audius} não aplicam esse middleware, pois a pesquisa do catálogo é pública na implementação observada.

\section{Integração Audius}
O serviço \arquivo{backend/src/modules/audius/audius.service.ts} consulta o gateway \url{https://api.audius.co}. A pesquisa utiliza \texttt{/v1/tracks/search}, enviando \texttt{query}, \texttt{app\_name}, \texttt{limit} e \texttt{offset}. A consulta individual utiliza \texttt{/v1/tracks/:id}.

A resposta externa é convertida para o modelo interno \texttt{Track}. \texttt{isPlayable} rejeita faixas sem identificador, indisponíveis ou com stream restrito. \texttt{normalizeTrack} converte título, artista, arte, duração e gênero para o formato interno.

O endereço de reprodução é produzido por \texttt{buildStreamUrl}:
\begin{lstlisting}
https://api.audius.co/v1/tracks/{trackId}/stream?app_name={appName}
\end{lstlisting}
O código comenta que esse endpoint estável pode responder com redirecionamento HTTP para o conteúdo de áudio. A URL é persistida na biblioteca como referência de reprodução, não como arquivo local.

A pesquisa calcula \texttt{hasMore} pela quantidade bruta de itens recebidos antes do filtro de faixas reproduzíveis. Assim, uma página filtrada vazia não é tratada automaticamente como fim do catálogo quando a API ainda retorna itens.

\section{Cliente HTTP e falhas}
A integração é isolada em \arquivo{audius-http.ts}, que utiliza Axios. Esse isolamento concentra configurações e tratamento da comunicação externa. A consulta individual lança erro 404 quando a resposta não contém identificador válido de faixa.

\section{Supabase Auth}
O Supabase Auth é utilizado para autenticação por e-mail e senha, conforme o README. No backend, o middleware lê \texttt{Authorization}, extrai o token no formato Bearer e chama \texttt{supabaseAuthClient.auth.getUser(accessToken)}. Token ausente, inválido ou expirado resulta em HTTP 401. Após validação, o identificador do usuário e o token são anexados à requisição.

\chapter{CORS e comunicação entre origens}
\section{Política implementada}
O servidor configura CORS por meio do pacote \texttt{cors}. A lista explícita de origens inclui \texttt{FRONTEND\_URL}, endereços locais e origens usadas pelas WebViews do Capacitor:
\begin{itemize}
\item \texttt{http://localhost:8100} e \texttt{http://127.0.0.1:8100};
\item \texttt{capacitor://localhost};
\item \texttt{http://localhost} e \texttt{https://localhost}.
\end{itemize}
Há também expressão regular para permitir endereços HTTP de rede privada nas faixas 10.x.x.x, 192.168.x.x e 172.16.x.x a 172.31.x.x, desde que utilizem a porta 8100. Requisições sem cabeçalho Origin também são aceitas. Origens não autorizadas recebem erro de CORS. A opção \texttt{credentials: true} habilita requisições com credenciais no contexto CORS, mantendo a validação da origem.

\section{O que CORS faz e o que não faz}
CORS é um mecanismo aplicado pelos navegadores para controlar se uma página pode ler respostas de outra origem. Não é autenticação, não criptografa dados e não torna confiável uma origem externa. A política CORS do backend controla o acesso do frontend à API Melodix; ela não concede, por si só, permissão para o navegador ler o áudio servido pelo Audius.

A leitura de amostras de áudio de outra origem depende de a resposta do recurso atender às regras de compartilhamento entre origens. Por isso, o cliente define \texttt{audio.crossOrigin = 'anonymous'} antes de atribuir a URL ao elemento de áudio. O servidor de mídia externo precisa fornecer cabeçalhos CORS adequados.

\section{Considerações sobre a lista de origens}
A expressão regular para redes privadas é uma conveniência para desenvolvimento e uso em rede local. Em implantação pública, convém avaliar se essa amplitude é necessária e restringir a lista às origens efetivamente utilizadas. CORS não substitui autenticação e autorização.

\chapter{Reprodução e streaming de áudio}
\section{Modelo de distribuição}
O Melodix não armazena os bytes das músicas no PostgreSQL. O backend consulta o Audius e devolve metadados e uma URL de stream. O navegador configura o elemento \texttt{HTMLAudioElement} com essa URL e solicita o conteúdo diretamente ao serviço externo.

Esse desenho reduz a responsabilidade de armazenamento e transferência de mídia do backend Melodix. A biblioteca persiste metadados suficientes para apresentar e localizar novamente uma faixa, além da URL de reprodução. A disponibilidade depende do serviço Audius e das condições de acesso ao stream.

\section{Sequência de reprodução}
\begin{enumerate}
\item O usuário seleciona uma faixa nos resultados.
\item O frontend recebe o objeto \texttt{Track}, incluindo \texttt{previewUrl}.
\item O \texttt{PlayerService} atualiza a faixa atual e zera o progresso quando há mudança de faixa.
\item O serviço prepara o grafo Web Audio na primeira reprodução, dentro de uma ação do usuário.
\item O elemento de áudio recebe a URL e executa \texttt{play()}.
\item Os eventos \texttt{waiting} e \texttt{playing} atualizam carregamento e reprodução.
\item \texttt{timeupdate} atualiza o progresso da interface.
\item \texttt{ended} desativa reprodução e carregamento e define o progresso como 100\%.
\end{enumerate}

\section{Controles e estados}
O serviço oferece reprodução, pausa, volume e busca por posição. O volume é recebido como percentual, limitado a 0--100 e convertido para 0--1, intervalo esperado pelo elemento de áudio. O progresso é expresso em percentual e convertido em segundos pela duração da faixa.

O player publica solicitações de faixa anterior e próxima. A fila é mantida pela lógica da tela de busca, que reage aos eventos e escolhe a faixa correspondente.

\section{Falhas de reprodução}
Quando uma faixa não possui \texttt{previewUrl}, o serviço limpa a reprodução e publica mensagem de indisponibilidade. A rejeição de \texttt{audio.play()} também é tratada. O grafo Web Audio é inicializado uma única vez por elemento de áudio. A proteção \texttt{audioGraphInitialized} evita chamadas repetidas a \texttt{createMediaElementSource} para o mesmo elemento.

\chapter{Visualização e análise de frequência}
\section{Grafo Web Audio}
O visualizador é implementado em \arquivo{src/app/components/audio-visualizer/audio-visualizer.component.ts}. O serviço de reprodução constrói o seguinte grafo:
\begin{lstlisting}
HTMLAudioElement
      |
      v
MediaElementAudioSourceNode
      |
      v
AnalyserNode
      |
      v
AudioContext.destination
\end{lstlisting}
O \texttt{AudioContext} representa o contexto de processamento. \texttt{MediaElementAudioSourceNode} conecta o elemento de áudio ao grafo. \texttt{AnalyserNode} obtém dados de frequência sem interromper a saída. A conexão com \texttt{destination} mantém o sinal encaminhado à saída de áudio.

\section{Parâmetros do analisador}
O código configura \texttt{analyserNode.fftSize = 64} e \texttt{smoothingTimeConstant = 0.8}. O tamanho FFT determina a resolução da análise. Para esse valor, \texttt{frequencyBinCount} contém 32 posições, metade de \texttt{fftSize}. O parâmetro de suavização reduz mudanças abruptas entre quadros consecutivos.

\section{Extração dos dados de frequência}
A cada quadro, o componente executa:
\begin{lstlisting}
analyser.getByteFrequencyData(data);
\end{lstlisting}
O método preenche um vetor de bytes com magnitudes de frequência no intervalo de 0 a 255. O componente normaliza cada posição dividindo por 255, obtendo intensidade entre 0 e 1.

O código utiliza até 85\% dos bins disponíveis, descartando a parte mais aguda do espectro. Cada barra seleciona um bin proporcional à sua posição e converte a magnitude em escala vertical:
\begin{lstlisting}
const value = data[bin] / 255;
const scale = 0.12 + value * 0.88;
element.style.transform = 'scaleY(' + scale.toFixed(3) + ')';
\end{lstlisting}
A opacidade também varia conforme a intensidade. A escala mínima mantém as barras visíveis em níveis baixos.

\section{Atualização em tempo real}
O ciclo é conduzido por \texttt{requestAnimationFrame}. \texttt{startLiveLoop} inicia o ciclo quando existe analisador, a faixa está efetivamente em reprodução e o carregamento terminou. \texttt{renderFrame} lê o espectro, atualiza os elementos e agenda o quadro seguinte.

As alterações diretas no DOM são realizadas com \texttt{NgZone.runOutsideAngular}. Isso evita executar a detecção de mudanças do Angular a cada quadro. O componente usa \texttt{ResizeObserver} para recalcular a quantidade de barras conforme a largura disponível, considerando largura e espaçamento.

\section{Detecção de silêncio e fallback}
O visualizador considera a possibilidade de o áudio tocar sem que o analisador receba amostras utilizáveis. Uma causa provável é a ausência de cabeçalhos CORS adequados no recurso externo. Nessa situação, o analisador pode fornecer valores zerados.

O componente soma os valores dos bins em cada quadro. Se ainda não observou sinal e acumula 120 quadros silenciosos, interrompe o modo real e retorna ao fallback CSS. A 60 quadros por segundo, esse limite representa aproximadamente dois segundos. O fallback não representa o espectro real da música; é uma animação independente para manter a interface funcional.

\section{Condições para análise real}
São necessárias: suporte do navegador à Web Audio API; criação do grafo em contexto permitido; carregamento do recurso com configuração CORS adequada; cabeçalhos do servidor de mídia que autorizem a origem; e reprodução efetiva, fora do buffering.

A configuração \texttt{crossOrigin = 'anonymous'} deve ocorrer antes de atribuir a URL. Ela não obriga o servidor remoto a permitir a origem. A autorização depende dos cabeçalhos retornados pelo Audius ou pelo host de mídia após redirecionamento.

\chapter{Banco de dados e persistência}
\section{Supabase e PostgreSQL}
O Supabase fornece autenticação e acesso ao PostgreSQL. A configuração local está em \arquivo{supabase/config.toml}. O projeto local é identificado como \texttt{spotify}; a configuração habilita API, banco, autenticação, Realtime e Storage. O limite de arquivo configurado para Storage local é 50 MiB.

A migração \arquivo{20260517192304\_auth\_owned\_library.sql} cria \texttt{playlists}, \texttt{playlist\_tracks} e \texttt{favorites}.

\section{Tabela playlists}
Armazena UUID, proprietário, nome, descrição, capa e datas de criação e atualização. \texttt{user\_id} referencia \texttt{auth.users(id)} com exclusão em cascata. O nome deve conter entre 1 e 80 caracteres após remoção de espaços externos. A descrição aceita até 240 caracteres.

\section{Tabela playlist\_tracks}
Associa faixas às playlists e armazena identificador externo, título, artista, álbum, miniatura, duração, URL de reprodução e posição. A combinação \texttt{playlist\_id, track\_id} é única. A chave estrangeira para playlist utiliza exclusão em cascata.

\section{Tabela favorites}
Registra favoritos por usuário. Armazena \texttt{user\_id}, \texttt{track\_id} e metadados suficientes para reconstruir o objeto de faixa. A combinação de usuário e faixa é única, permitindo operação \texttt{upsert}.

\section{Índices e integridade}
A migração cria índices para consultas ordenadas por usuário e data em playlists e favoritos, além de índice para recuperar faixas de uma playlist pela ordem. Estabelece restrições de duração não negativa, limites textuais e unicidade.

\section{Dados persistidos e não persistidos}
O banco armazena informações de organização e metadados de faixas, incluindo identificador Audius e URL de reprodução. As tabelas da biblioteca não armazenam o conteúdo binário do áudio. O arquivo musical permanece sob responsabilidade da infraestrutura externa. A configuração Supabase habilita Storage, mas isso não comprova que a aplicação esteja armazenando músicas nesse serviço. A implementação analisada utiliza tabelas PostgreSQL para metadados.

\chapter{Segurança no backend e no banco}
\section{Autenticação Bearer}
O middleware \arquivo{backend/src/modules/auth/auth.middleware.ts} exige o cabeçalho \texttt{Authorization} no formato Bearer. O token é validado com \texttt{auth.getUser(accessToken)}. Tokens ausentes ou inválidos resultam em HTTP 401. Após validação, o identificador do usuário e o token são anexados à requisição.

\section{Row Level Security}
A migração ativa RLS nas três tabelas da biblioteca. As políticas permitem operações para usuários autenticados e verificam a correspondência entre a identidade da sessão e o proprietário do registro.

Na tabela \texttt{playlists}, as políticas comparam \texttt{auth.uid()} com \texttt{user\_id}. A inserção utiliza \texttt{WITH CHECK}. Em \texttt{playlist\_tracks}, as políticas verificam a propriedade por meio de consulta à tabela \texttt{playlists}. Em \texttt{favorites}, as políticas verificam diretamente \texttt{user\_id}. Leitura, inserção, atualização e exclusão são autorizadas separadamente.

RLS atua no banco: mesmo que uma chamada alcance a API de dados, as políticas limitam os registros acessíveis à identidade autenticada.

\section{Separação de chaves Supabase}
O arquivo \arquivo{backend/src/lib/supabase.ts} instancia clientes com a chave publicável e também um cliente administrativo com \texttt{SUPABASE\_SERVICE\_ROLE\_KEY}. Define ainda \texttt{createSupabaseUserClient(accessToken)}, que injeta o token do usuário no cabeçalho Authorization.

As operações de \arquivo{library.service.ts} utilizam \texttt{createSupabaseUserClient}. Portanto, a biblioteca opera no contexto do usuário e pode ser submetida às políticas RLS. O cliente administrativo está definido, mas sua declaração não demonstra uso nas operações regulares da biblioteca.

A chave de serviço possui privilégios elevados e deve permanecer exclusivamente no servidor. Não deve ser incluída em arquivos de ambiente do frontend, código entregue ao navegador, repositório público ou logs.

\section{Validação e limites de requisição}
Zod valida variáveis de ambiente na inicialização. O backend possui módulo de esquemas para dados da biblioteca. A validação de negócio deve ocorrer no servidor, sem depender dos controles visuais do cliente. O parser JSON do Express está limitado a 1 MB, reduzindo a aceitação de corpos excessivamente grandes.

\section{CORS e limites de segurança}
CORS restringe quais origens de navegador podem ler respostas do backend, mas não substitui a validação de token. Clientes que não são navegadores podem emitir requisições HTTP independentemente da política CORS do navegador. A autorização das rotas privadas depende do middleware Bearer e das políticas RLS.

\chapter{Serviços da biblioteca e operações de dados}
\section{Listagem de playlists}
\texttt{listPlaylists} cria cliente Supabase associado ao token, consulta playlists ordenadas pela data de criação e recupera as faixas das playlists retornadas. \texttt{mapPlaylist} agrupa faixas pelo identificador da playlist e converte registros para o modelo de domínio.

Essa estratégia usa uma consulta para playlists e outra para as faixas relacionadas, evitando uma consulta individual para cada playlist.

\section{Criação, atualização e exclusão}
A criação insere o identificador do usuário, nome e descrição. Depois consulta novamente a playlist para retornar a representação completa. A atualização altera nome, descrição e data de atualização. A exclusão remove o registro e verifica se alguma linha foi afetada; na ausência, retorna 404.

\section{Inclusão e remoção de faixas}
Na inclusão, o serviço consulta a playlist e verifica se a faixa já existe. Se existir, atualiza metadados. Caso contrário, consulta a maior posição atual, calcula a próxima posição e insere o registro. A remoção verifica a existência e propriedade da playlist e exclui a associação correspondente.

\section{Favoritos}
A inclusão utiliza \texttt{upsert} com conflito em \texttt{user\_id,track\_id}, assegurando uma chave única por usuário e faixa. A remoção filtra pelo identificador da faixa; RLS limita o efeito aos registros que o usuário pode excluir.

\chapter{Fluxos principais}
\section{Pesquisa de músicas}
\begin{enumerate}
\item A tela recebe o texto digitado.
\item O serviço Angular envia consulta, limite e deslocamento ao backend.
\item O controlador Audius encaminha os parâmetros ao serviço.
\item O serviço consulta o endpoint Audius com Axios.
\item Os resultados são filtrados por disponibilidade e restrição de stream.
\item Os objetos externos são normalizados para o tipo interno \texttt{Track}.
\item O backend retorna faixas e \texttt{hasMore}.
\item A tela atualiza a lista e solicita páginas adicionais durante a rolagem.
\end{enumerate}

\section{Acesso à biblioteca}
\begin{enumerate}
\item O frontend envia o token no cabeçalho Authorization.
\item O middleware valida o token com Supabase Auth.
\item O identificador do usuário é associado à requisição.
\item O serviço cria cliente Supabase com o token daquele usuário.
\item A consulta é executada no PostgreSQL e submetida às políticas RLS.
\item Os registros são convertidos para modelos de playlist e faixa.
\item O backend devolve a resposta ao frontend.
\end{enumerate}

\section{Reprodução e visualizador}
A reprodução utiliza a URL de stream da faixa selecionada. O PlayerService controla o elemento de áudio e emite estados à interface. O visualizador observa o analisador e os estados de reprodução e carregamento. Quando o áudio está disponível, o componente lê os bins de frequência a cada quadro. Sem sinal válido, retorna à animação CSS.

\chapter{Pontos de atenção e recomendações}
\section{Gestão de segredos}
O arquivo \arquivo{backend/.env.example} documenta variáveis necessárias. O arquivo real \arquivo{backend/.env} deve permanecer fora do versionamento. Recomenda-se verificar \arquivo{.gitignore} e o histórico do repositório para assegurar que chaves de serviço não tenham sido publicadas.

\section{Uso do cliente administrativo}
O cliente \texttt{supabaseAdminClient} deve ser utilizado somente quando uma operação privilegiada for realmente necessária. Para operações associadas a um usuário, a abordagem atual de criar cliente com o token preserva a avaliação de RLS.

\section{Política CORS}
As origens devem refletir os ambientes suportados. A expressão regular de rede privada é útil para desenvolvimento local, mas deve ser avaliada antes de ser mantida em produção.

\section{Dependência do Audius}
Pesquisa e reprodução dependem da disponibilidade, políticas e endpoints do Audius. Uma indisponibilidade externa pode afetar catálogo ou reprodução sem que o backend Melodix esteja fora do ar. O tratamento de erros deve distinguir falhas de catálogo, de stream e internas.

\section{URL de reprodução}
O projeto persiste a URL de stream junto aos metadados. A reprodução futura depende da estabilidade e das regras do endpoint Audius. O código constrói uma URL de stream com \texttt{app\_name}, em vez de armazenar URL temporária assinada.

\section{Observabilidade e testes}
O backend possui \texttt{/health}. Recomenda-se complementar essa verificação com logs estruturados, métricas de latência, acompanhamento de erros externos e testes de autenticação, RLS, paginação, normalização de faixas e comportamento do player.

\chapter{Inventário de arquivos relevantes}
\begin{longtable}{p{0.49\textwidth}p{0.42\textwidth}}
\toprule
\textbf{Arquivo ou diretório} & \textbf{Responsabilidade}\\
\midrule
\arquivo{README.md} & Funcionalidades, arquitetura e execução.\\
\arquivo{package.json} & Dependências e scripts do frontend.\\
\arquivo{backend/package.json} & Dependências e scripts do backend.\\
\arquivo{src/app/core/core.module.ts} & Serviços centrais e interceptor.\\
\arquivo{src/app/app.component.ts} & Navegação, mini-player e mensagens.\\
\arquivo{src/app/core/services/player.service.ts} & Reprodução, estados e grafo Web Audio.\\
\arquivo{src/app/core/services/audius.service.ts} & Cliente frontend da API Audius interna.\\
\arquivo{src/app/components/player/player.component.ts} & Controles e interações do player.\\
\arquivo{src/app/components/audio-visualizer/audio-visualizer.component.ts} & Leitura de frequência e barras.\\
\arquivo{src/app/pages/search/} & Pesquisa, resultados, paginação e seleção.\\
\arquivo{backend/src/server.ts} & Inicialização Express, CORS e rotas.\\
\arquivo{backend/src/config/env.ts} & Validação de configuração com Zod.\\
\arquivo{backend/src/lib/supabase.ts} & Clientes Supabase e credenciais.\\
\arquivo{backend/src/modules/auth/auth.middleware.ts} & Validação Bearer.\\
\arquivo{backend/src/modules/audius/} & Rotas, controladores, cliente HTTP e serviço Audius.\\
\arquivo{backend/src/modules/library/} & Operações de playlists e favoritos.\\
\arquivo{supabase/config.toml} & Configuração local do Supabase.\\
\arquivo{supabase/migrations/20260517192304\_auth\_owned\_library.sql} & Tabelas, índices, RLS e permissões.\\
\arquivo{supabase/migrations/20260518002513\_spotify\_connections.sql} & Tabela de conexões Spotify.\\
\arquivo{capacitor.config.ts} & Configuração do Capacitor.\\
\bottomrule
\end{longtable}

\chapter{Conclusão}
A arquitetura do Melodix separa interface, regras de backend, catálogo externo e persistência da biblioteca. Angular, Ionic e RxJS estruturam a experiência cliente. Express e TypeScript fornecem a API intermediária. Audius fornece dados do catálogo e infraestrutura de streaming. Supabase Auth e PostgreSQL sustentam a identidade e os dados pessoais da biblioteca.

Os mecanismos de segurança observados incluem validação de tokens Bearer, clientes Supabase vinculados ao token do usuário, políticas Row Level Security, validação de configuração por Zod, limite do corpo JSON e política CORS por origem. A chave de serviço Supabase é uma credencial privilegiada e deve permanecer restrita ao servidor.

A reprodução utiliza o elemento de áudio do navegador. A análise de frequência é feita pela Web Audio API por meio de um grafo com \texttt{MediaElementAudioSourceNode} e \texttt{AnalyserNode}. A configuração CORS do elemento e os cabeçalhos retornados pelo servidor de mídia determinam se o analisador terá acesso às amostras. O fallback CSS mantém a apresentação visual quando a análise real não é possível.

Esta documentação descreve o comportamento identificado nos arquivos examinados. Configurações de implantação, disponibilidade de serviços externos e versões efetivamente instaladas podem variar entre ambientes.

\appendix
\chapter{Referências técnicas}
\begin{itemize}
\item Angular: \url{https://angular.dev/}
\item Ionic Framework: \url{https://ionicframework.com/docs}
\item RxJS: \url{https://rxjs.dev/}
\item Capacitor: \url{https://capacitorjs.com/docs}
\item Node.js: \url{https://nodejs.org/docs/}
\item Express: \url{https://expressjs.com/}
\item TypeScript: \url{https://www.typescriptlang.org/docs/}
\item Zod: \url{https://zod.dev/}
\item Supabase Auth: \url{https://supabase.com/docs/guides/auth}
\item Supabase RLS: \url{https://supabase.com/docs/guides/database/postgres/row-level-security}
\item Audius API: \url{https://docs.audius.org/}
\item MDN Web Audio API: \url{https://developer.mozilla.org/docs/Web/API/Web_Audio_API}
\item MDN CORS: \url{https://developer.mozilla.org/docs/Web/HTTP/CORS}
\end{itemize}
\end{document}
